\chapter{Introduction to Computer Algebra Systems}\label{cas}

A computer algebra system is a programming language that supports symbolic computation.
Computer algebra systems treat unbound variables as symbols.
Using a computer algebra system, one can perform computations such as $x + x \rightarrow 2x$ and $(x + y)^2 \rightarrow x^2 + 2 x y + y^2$.
Egison also supports such symbolic computation.
This chapter explores what kinds of computations are possible with a computer algebra system.

\section{Declaring Symbols}\label{cas-symbol}

A computer algebra system treats unbound variables as symbols.
In Egison, symbols must be declared using the \lstinline{declare symbol} statement before use.
Types can be specified for symbols; if omitted, the \lstinline{MathExpr} type is used by default.
\begin{lstlisting}[language=egison,numbers=none]
declare symbol x, y, z : MathExpr

x -- x
\end{lstlisting}

Egison, as a computer algebra system, has built-in definitions for addition and multiplication of symbols.
Therefore, computations such as $x + x \rightarrow 2x$ and $(x + y)^2 \rightarrow x^2 + 2 x y + y^2$ are possible.
\begin{lstlisting}[language=egison,numbers=none]
x + x -- 2x
(x + y)^2 -- x^2 + 2 * x * y + y^2
\end{lstlisting}

Egison automatically expands mathematical expressions into sum-of-products normal form.
Sum-of-products normal form is a form of expression where multiplication appears on the inside and addition on the outside.
Therefore, $(x + y)^2$ is expanded to $x^2 + 2 x y + y^2$.

The ordering of symbol factors in the output expression is implemented to follow the order in which they appear in the user's input expression.
For example, if we swap $x$ and $y$ in $(x + y)^2$ above and compute $(y + x)^2$, the second term becomes $2 y x$.
\begin{lstlisting}[language=egison,numbers=none]
(y + x)^2 -- y^2 + 2 * y * x + x^2
\end{lstlisting}

\section{Symbolic Function Application}\label{cas-symbol2}

For functions such as \lstinline{sqrt x} and \lstinline{exp x} that may take symbols as arguments, the library defines them to halt evaluation and return the expression as a symbolic value.
\begin{lstlisting}[language=egison,numbers=none]
declare symbol x : MathExpr

sqrt 4 -- 2
sqrt x -- sqrt x
exp 1 -- e
exp x -- exp x
\end{lstlisting}

For function applications that cannot be further evaluated, such as \lstinline{sqrt 2}, the library also defines them to halt evaluation and return the expression as is.
Simplifications such as $\sqrt{8} \rightarrow 2 \sqrt{2}$ are also defined in the library.
\begin{lstlisting}[language=egison,numbers=none]
sqrt 2 -- sqrt 2
sqrt 8 -- 2 * sqrt 2
\end{lstlisting}
These definitions use single quotes, which are introduced in Section~\ref{mexpr-squote}.


\section{Simplification of Expressions}\label{cas-simplify}

Simplifications that one performs when working with mathematical expressions on paper, such as $\sqrt{x} \cdot \sqrt{x} = x$ and $sin^2\theta + cos^2\theta = 1$, are implemented in many computer algebra systems including Egison.
\begin{lstlisting}[language=egison,numbers=none]
declare symbol x, `$\theta$` : MathExpr

sqrt x * sqrt x -- x
(sin `$\theta$`)^2 + (cos `$\theta$`)^2 -- 1
\end{lstlisting}
In Egison, such simplification functionality is built into the language, implemented in Haskell using Sweet Egison (a Haskell library that provides Egison's pattern matching capabilities).

In addition to the above, simplification rules for complex numbers and other entities are also implemented.
\begin{lstlisting}[language=egison,numbers=none]
i^2 -- -1
w^3 -- 1  
w + w^2 -- -1  
\end{lstlisting}
The simplification rules are defined at \url{https://github.com/egison/egison/blob/master/hs-src/Language/Egison/Math/Normalize.hs}.

Automating the simplification of mathematical expressions is extremely difficult, and Egison provides only simple algorithms.
For example, algorithms based on Gr\"obner basis theory are useful for polynomial division, but they are not implemented in Egison.
As a result, Egison may not be able to properly simplify complex polynomial fractions.
For univariate polynomials, however, reduction by the Euclidean algorithm, such as $(x^2-1)/(x-1) \rightarrow x+1$, is implemented.

Wolfram Language (Mathematica) is capable of such advanced expression simplification.
Running the Egison interpreter with the ``\texttt{-M mathematica}'' option outputs expressions in a format readable by the Wolfram Language.
By feeding this output to the Wolfram Language interpreter, one can leverage its advanced simplification capabilities.
\begin{lstlisting}[language=egison,numbers=none]
$ egison -M mathematica
> x + sqrt 2
#mathematica|x + Sqrt[2]|#
\end{lstlisting}

\section{Declaring Simplification Rules}\label{cas-rule}

In addition to the built-in simplification rules introduced in the previous section, users can declare their own simplification rules with the \lstinline{declare rule} statement.
For example, if we introduce a symbol \lstinline{sqrt2} representing $\sqrt{2}$ and register the rule $\sqrt{2}^2 = 2$, the rule is applied automatically during every simplification.
\begin{lstlisting}[language=egison,numbers=none]
declare symbol sqrt2
declare rule auto term sqrt2^2 = 2

(1 + sqrt2)^2 -- 2 * sqrt2 + 3
\end{lstlisting}
Mathematically, this operation introduces the algebraic extension $\mathbb{Z}[\sqrt{2}]$, i.e., the quotient $\mathbb{Z}[X]/(X^2-2)$ of a polynomial ring.
Simplification rules in Egison are global, but since symbols are globally unique, a rule can only act on expressions in which its symbol occurs.
In other words, the uniqueness of symbols plays the role of scoping.
The Gaussian integers $\mathbb{Z}[i]$ (with $i^2 = -1$), the dual numbers (with $\varepsilon^2 = 0$), and roots of unity are algebraic extensions introduced by the same pattern.

\section{Selecting Canonical Forms with Types}\label{cas-tower}

The same mathematical expression can be represented by several different data structures.
For example,
\[ (2+3i) + (1-i)x + 4x^2 \]
can be held as a flat polynomial over the atoms $\{i, x\}$, or as a polynomial in $x$ whose coefficients are Gaussian integers in $\mathbb{Z}[i]$.
Which canonical form is appropriate depends on what one intends to do next (differentiate with respect to $x$, or exploit properties of $\mathbb{Z}[i]$), so it cannot be determined from the value alone.
In Egison, \textbf{type annotations select the canonical form}.
\begin{lstlisting}[language=egison,numbers=none]
declare symbol i, x

def v := (2 + 3 * i) + (1 - i) * x + 4 * x^2

def flat : Poly Integer [i, x] := v
inspect flat
-- "4 * x^2 - i x + x + 3 * i + 2 : Poly Integer [i, x]"

def byX : Poly (Poly Integer [i]) [x] := v
inspect byX
-- "4 * x^2 + (- i + 1) * x + 3 * i + 2 : Poly (Poly Integer [i]) [x]"
\end{lstlisting}
The type \lstinline{Poly Integer [i, x]} denotes integer-coefficient polynomials over the atom set $\{i, x\}$.
Nesting a polynomial type in the coefficient position, as in \lstinline{Poly (Poly Integer [i]) [x]}, denotes the canonical form ``polynomials in $x$ with coefficients in $\mathbb{Z}[i]$''.
A type annotation does more than classify the value: it reorganizes the value's internal structure into the shape of the type (reshape).
Reshaping never changes the value.
Indeed, the difference of the two canonical forms computes to $0$.
\begin{lstlisting}[language=egison,numbers=none]
flat - byX -- 0
\end{lstlisting}
Arithmetic between differently shaped representations simplifies correctly because of the convention that \textbf{operations always return results in the default flat form}.
Nested canonical forms appear exactly where annotations request them and never propagate silently across operations.

The \lstinline{inspect} used above displays a value together with the type read off the value (its observed type).
Frequently used types can be given names with the \lstinline{declare cas-type} statement.
Aliases are transparent: displayed types use the expanded original form.
\begin{lstlisting}[language=egison,numbers=none]
declare cas-type GaussianPolyX := Poly (Poly Integer [i]) [x]

def byX2 : GaussianPolyX := v
inspect byX2
-- "4 * x^2 + (- i + 1) * x + 3 * i + 2 : Poly (Poly Integer [i]) [x]"
\end{lstlisting}

\section{Subtyping and Implicit Promotion}\label{cas-subtype}

The CAS types are ordered by value-preserving embeddings.
For example, integers embed into rationals (\lstinline{Integer} $\subset$ \lstinline{Frac Integer}), and rationals embed into rational-coefficient polynomials as constants.
When values of different types are mixed in one operation, they are automatically promoted to the least type containing both (their join).
\begin{lstlisting}[language=egison,numbers=none]
def p : Poly Integer [x] := x + 1
def q : Frac Integer := 1 / 2

def r := p + q
inspect r -- "x + 3 / 2 : Poly (Frac Integer) [x]"
\end{lstlisting}
When the atom sets differ, values are promoted to the type over the union of the atom sets.
\begin{lstlisting}[language=egison,numbers=none]
declare cas-type GaussianInt := Poly Integer [i]
declare cas-type Zsqrt2 := Poly Integer [sqrt2]

def a : GaussianInt := 1 + i
def b : Zsqrt2 := 1 + sqrt2

def c := a + b
inspect c -- "sqrt2 + i + 2 : Poly Integer [i, sqrt2]"
\end{lstlisting}
The sum of a $\mathbb{Z}[i]$ value and a $\mathbb{Z}[\sqrt{2}]$ value lands in $\mathbb{Z}[i, \sqrt{2}]$ without any annotation.

The built-in order deliberately does not relate the flat form (\lstinline{Poly Integer [i, x]}) and the nested form (\lstinline{Poly (Poly Integer [i]) [x]}) of the previous section.
They are different canonical forms of the same set of values, and which one to promote to is the user's choice.
This relation can be declared by the user with the \lstinline{declare cas-subtype} statement.
\begin{lstlisting}[language=egison,numbers=none]
declare cas-subtype Poly Integer [i, x] <: Poly (Poly Integer [i]) [x]
\end{lstlisting}
Each declared edge is checked to keep the join of every pair of types unique over the whole order (the join-semilattice property).
A declaration that would break uniqueness is rejected together with a suggestion of the edge to declare first, so the order always stays sound.

\section{Quotient Types}\label{cas-quotient}

Quotients such as $\mathbb{Z}/7\mathbb{Z}$ (the world of remainders modulo 7) behave differently from the types so far.
$12$ and $5$ are different values in $\mathbb{Z}$ but equal in $\mathbb{Z}/7\mathbb{Z}$.
Since \textbf{equality itself depends on the type}, quotients cannot join the order of value-preserving embeddings of the previous section.
In Egison, quotient types are declared outside the order with the \lstinline{declare cas-quotient} statement.
\begin{lstlisting}[language=egison,numbers=none]
declare cas-quotient Mod7 := Integer by (\n -> modulo n 7)

def a : Mod7 := projMod7 5
def b : Mod7 := projMod7 4

a + b -- 2
projMod7 12 == projMod7 5 -- True
\end{lstlisting}
The declaration automatically derives the projection \lstinline{projMod7} from integers into \lstinline{Mod7}, the representative extraction \lstinline{reprMod7}, the ring operations, and the equality test.
Every operation reduces its representative, so values always stay in the range $0$ to $6$.
At declaration time, the system also tests, on sample points, that the reduction function is idempotent and compatible with addition and multiplication (a congruence).

Mixing a \lstinline{Mod7} value with a raw integer (as in \lstinline{a + 10}) is a type error.
This prevents silently confusing equality in $\mathbb{Z}$ with equality in $\mathbb{Z}/7\mathbb{Z}$; crossing between the two worlds is explicit, via \lstinline{projMod7} and \lstinline{reprMod7}.
Functions can also be defined directly on the quotient.
For example, the inverse in $\mathbb{F}_7$ is one line by Fermat's little theorem ($z^{-1} = z^5$).
\begin{lstlisting}[language=egison,numbers=none]
def inv7 (z : Mod7) : Mod7 := z * z * z * z * z

reprMod7 (inv7 (projMod7 3)) -- 5
\end{lstlisting}

\section{A Program to Solve Quadratic Equations -- Writing Algorithms that Process Expressions}\label{cas-quadratic}

In this section, we will write an algorithm to solve quadratic equations as a program.
We define a function \lstinline{qF} that takes a quadratic equation and returns its solutions, as shown below.
\begin{lstlisting}[language=egison,numbers=none]
qF (x ^ 2 + x + 1) x
-- ((-1 + i * sqrt 3) / 2, (-1 - i * sqrt 3) / 2)

qF (x ^ 2 + b * x + c) x
-- ((- b + sqrt (b^2 - 4 * c)) / 2, (- b - sqrt (b^2 - 4 * c)) / 2)

qF (a * x ^ 2 + b * x + c) x
-- ((- b + sqrt (b^2 - 4 * a * c)) / (2 * a), (- b - sqrt (b^2 - 4 * a * c)) / (2 * a))

qF (a * x ^ 2 + 2 * b * x + c) x
-- ((- b + sqrt (b^2 - a * c)) / a, (- b - sqrt (b^2 - a * c)) / a)
\end{lstlisting}
\lstinline{qF} can be defined by simply writing out the quadratic formula $\frac{-b \pm \sqrt{b^2 - 4 a c}}{2 a}$ directly.
The following program defines \lstinline{qF} using an auxiliary function \lstinline{qF'}.
\lstinline{qF'} takes the quadratic coefficient as the first argument, the linear coefficient as the second, and the constant term as the third.
For example, \lstinline{qF' 1 0 2} returns the solutions of the quadratic equation $x^2 + 2 = 0$.
\lstinline{coefficients} is a function that takes a polynomial and a symbol, and returns the list of coefficients with respect to that symbol.
\lstinline{qF'} is simply defined as a function that returns a tuple consisting of $\frac{-b + \sqrt{b^2 - 4 a c}}{2 a}$ and $\frac{-b - \sqrt{b^2 - 4 a c}}{2 a}$.
\begin{lstlisting}[language=egison,numbers=none]
def qF f x :=
  match coefficients f x as list mathExpr with
    | [$a_0, $a_1, $a_2] -> qF' a_2 a_1 a_0

def qF' a b c :=
  ( ((- b) + sqrt (b ^ 2 - 4 * a * c)) / 2 * a
  , ((- b) - sqrt (b ^ 2 - 4 * a * c)) / 2 * a )
\end{lstlisting}
However, since the above formulation lacks an algorithmic feel, we will instead write an algorithm that derives the quadratic formula through completing the square.
Completing the square is the operation of transforming a quadratic expression such as $x^2 + b x + c$ into the form $(x + b / 2)^2 - b^2 / 4 + c$.
\begin{lstlisting}[language=egison]
def qF' a b c :=
  match (a, b, c) as (mathExpr, mathExpr, mathExpr) with
    | (#1, #0, _) -> (sqrt (- c), - sqrt (- c))
    | (#1, _, _) ->
      (2)#((- (b / 2)) + %1, (- (b / 2)) + %2)
        (withSymbols [x, y]
          qF (substitute [(x, y - b / 2)] (x ^ 2 + b * x + c)) y)
    | (_, _, _) -> qF' 1 (b / a) (c / a)
\end{lstlisting}
When describing the algorithm for solving quadratic equations using completing the square, \lstinline{qF'} becomes a function with three match clauses.
The match clause on line 3 handles quadratic equations of the form $x^2 + c = 0$.
This clause returns the result $\sqrt{c}$ and $-\sqrt{c}$.
The match clause on lines 4--7 handles quadratic equations of the form $x^2 + b x + c = 0$, where the leading coefficient is $1$.
This clause is the core of the function, performing the completing-the-square operation.
\lstinline{substitute} is an Egison library function.
\lstinline{substitute} applies the substitution specified in the first argument to the expression in the second argument.
Here, the substitution \lstinline{(x, y - b / 2)} is specified in the first argument, so the substitution $x = y - b / 2$ is applied to the expression $x^2 + b x + c$.
The first argument of \lstinline{substitute} takes a list, allowing multiple simultaneous substitutions.
After substituting $x = y - b / 2$, this equation is transformed into a quadratic equation in $y$ that matches the clause on line 3.
Therefore, by passing this transformed equation to the \lstinline{qF} function, we can obtain the solutions for $y$.
By adding $\frac{-b}{2}$ to these solutions for $y$, we obtain the solutions for the original variable $x$.
The \lstinline{withSymbols} expression is a syntax for generating local symbols.
The symbols specified in the list of the first argument (in this case, \lstinline{x} and \lstinline{y}) can be used as symbols within the expression of the second argument.
Thanks to this feature of \lstinline{withSymbols}, even if \lstinline{x} or \lstinline{y} are bound to specific integers elsewhere in the program, we do not need to worry about that inside the \lstinline{withSymbols} expression.
The match clause on line 8 handles quadratic equations where the leading coefficient is not $1$.
It divides the equation by the leading coefficient to convert it to an equation with leading coefficient $1$, and passes it to the \lstinline{qF} function again.

Although more complex than quadratic equations, it is interesting to write programs that solve cubic and quartic equations.
These programs are available in the \texttt{sample/math/algebra} directory of the Egison source code.

\section{A Program for Differentiation -- Pattern Matching on Expressions}\label{cas-derivative}

Egison has a library function \lstinline{d/d} for performing differentiation.
This function returns the result of differentiating the first argument with respect to the symbol given as the second argument.
By naming the function \lstinline{d/d}, differentiation can be expressed in a form close to mathematical notation.
Egison allows \lstinline{/} to be used in variable names.
\begin{lstlisting}[language=egison,numbers=none]
d/d x x -- 1
d/d (x^2) x -- 2 * x
d/d (exp x) x -- exp x
d/d (log x) x -- 1 / x
d/d (x * log x) x -- log x + 1
d/d (1 / log x) x -- -1 / (x * (log x)^2)
\end{lstlisting}
This section presents the implementation of this function.

\lstinline{d/d} is defined as follows.
\begin{lstlisting}[language=egison]
def d/d f x :=
  match f as mathExpr with
    -- Differentiation of symbols
    | #x -> 1
    | ?isSymbol -> 0
    -- Differentiation of function applications
    | #'exp $g -> exp g * d/d g x
    | #'log $g -> 1 / g * d/d g x
    | #'sqrt $g -> 1 / (2 * sqrt g) * d/d g x
    | #'(^) $g $h -> f * d/d (log g * h) x
    | #'cos $g -> (- sin g) * d/d g x
    | #'sin $g -> cos g * d/d g x
    -- Differentiation of constants
    | #0 -> 0
    | _ * #1 -> 0
    -- Differentiation of terms
    | #1 * $fx ^ $n -> n * fx ^ (n - 1) * d/d fx x
    | $a * $fx ^ $n * $r -> a * d/d (fx ^' n) x * r + a * fx ^' n * d/d r x
    -- Differentiation of polynomials
    | poly $ts -> sum (map 1#(d/d %1 x) ts)
    -- Differentiation of quotients
    | $p1 / $p2 ->
      let p1' := d/d p1 x
          p2' := d/d p2 x
       in (p1' * p2 - p2' * p1) / p2 ^ 2
\end{lstlisting}
The first argument $f$ is pattern-matched as a mathematical expression.
Lines 4--5 define the differentiation of symbols.
If \lstinline{f} is \lstinline{x} itself, it returns $1$; if \lstinline{f} is a symbol other than \lstinline{x}, it returns $0$.
Lines 7--12 describe the chain rule $\frac{df(g(x))}{dx} = \frac{df(g(x))}{dg(x)} \frac{dg(x)}{x}$ for several basic function applications.
For example, line 7 handles the differentiation of the exponential function.
The body of this match clause represents $\frac{de^{g(x)}}{dx} = e^{g(x)} \frac{dg(x)}{dx}$.
Line 8 handles the differentiation of the logarithmic function.
The body of this match clause represents $\frac{dlog(g(x))}{dx} = \frac{1}{g(x)} \frac{dg(x)}{dx}$.
Lines 14--15 define differentiation for constant terms.
Constant terms differentiate to $0$.
Lines 17--18 define differentiation for terms that are products of the previously defined differentiable terms.
This uses the facts that $x^n$ differentiates to $n x^{n-1}$ and that $f g$ differentiates to $f' g + f g'$.
The chain rule is also used here.
Line 20 defines the differentiation of polynomials.
Since differentiation of monomials has already been defined, for a sum of monomials, we simply differentiate each term and sum the results.
Lines 22--24 define differentiation for expressions with polynomials in both the numerator and denominator.
This is defined using the quotient rule for differentiation.

\section{Computation with Vectors and Matrices -- Index Notation}\label{cas-linear}

A distinctive feature of Egison as a computer algebra system is the ability to concisely describe computations involving vectors, matrices, and their generalization, tensors.
This is because Egison supports tensor index notation as a programming language feature -- a notation invented in the early twentieth century during research in differential geometry and general relativity.

Index notation expresses various kinds of tensor multiplication by attaching indices to tensors.
For example, there are three types of multiplication between vectors: the tensor product, the Hadamard product, and the inner product. These are expressed by using different index patterns as follows.
\begin{lstlisting}[language=egison,numbers=none]
[| x1, x2 |]_i .[| y1, y2 |]_j -- [| [| x1 * y1, x1 * y2 |], [| x2 * y1, x2 * y2 |] |]_i_j
[| x1, x2 |]_i .[| y1, y2 |]_i -- [| x1 * y1, x2 * y2 |]_i
[| x1, x2 |]~i .[| y1, y2 |]_i -- x1 * y1 + x2 * y2
\end{lstlisting}
As shown above, in Egison, vectors are represented by enclosing their components in \lstinline{[| |]}.
Matrices are represented as vectors of vectors.
From a functional programming perspective, index notation can be explained as adding indices as extra parameters to tensor multiplication functions, allowing a single function to serve multiple roles.

Other programming languages also have index notation as a language feature.
Among them, Egison's distinguishing characteristic is that it incorporates index notation in a way that harmonizes well with functional programming.
Thanks to this, higher-order functions and index notation can be used in combination.
For example, an expression like $g_{i_{1}j_{1}} g_{i_{2}j_{2}} ... g_{i_{n}j_{n}}$ can be written using the \lstinline{foldl} function as follows.
\begin{lstlisting}[language=egison,numbers=none]
foldl . 1 (map (\x -> g_[i_x]_[j_x]) [1..n])
\end{lstlisting}

Since index notation is an important language feature of Egison, it is covered in detail in Chapter~\ref{tensor} and Chapter~\ref{diffgeo}.

\section{Output Format for Expressions}

The output format can be specified using the ``\texttt{-M}'' option.
Running ``\texttt{egison -M latex}'' outputs results in LaTeX format.
\begin{lstlisting}[language=egison,numbers=none]
$ egison -M latex
Egison Version 3.7.14 (C) 2011-2018 Satoshi Egi
https://www.egison.org
Welcome to Egison Interpreter!
> (+ (sqrt a) b)
#latex|\sqrt{a} + b|#
\end{lstlisting}
Running ``\texttt{egison -M mathematica}'' outputs results in Mathematica format.
\begin{lstlisting}[language=egison,numbers=none]
$ egison -M mathematica
Egison Version 3.7.14 (C) 2011-2018 Satoshi Egi
https://www.egison.org
Welcome to Egison Interpreter!
> (+ (sqrt a) b)
#mathematica|Sqrt[a] + b|#
\end{lstlisting}

The LaTeX output also integrates with Jupyter Notebook.
By using the Egison plugin for Jupyter Notebook, computation results are displayed interactively as mathematical expressions, as shown in Figures~\ref{fig:riemann-curavture-of-S2} and~\ref{fig:hodge-laplacian}.

\begin{figure}[t]
  \centering
  \includegraphics[width=14.0cm,bb=0 0 1618 870]{jupyter-sample.png}
  \caption{Running Egison on Jupyter Notebook}
  \label{fig:cas-jupyter}
\end{figure}
